\subsection{Test vehicle verification and testing after manufacturing}
\label{sec:test-vehicle-testing}

The testplan for the testchip is constituted of 7 steps, described in Fig. \ref{fig:verif-plan}.
The verification of the on-board voltage regulation is checked with a voltmeter, to ensure that a regulated 12V is well established.
The on-chip regulation is checked with a voltmeter and observed with an oscilloscope to validate the behavior.
All regulators from both boards and inside the testchip were found to operate as expected.

\begin{figure}[!h]
  \centering
  \includegraphics[width=0.7\textwidth]{src/3/figures/verification_plan.pdf}
  \caption{Test plan}
  \label{fig:verif-plan}
\end{figure}


\subsubsection{Optical communication}

% What is this step
The communication between the optical transceiver is verified.
Both the microcontroller and testchip were disconnected prior to injecting test frames on each transceiver.
Then, the output of the matching receiver was observed to ensure that the frame was properly transmitted.
Overall, minor issues were identified and corrected, and the boards were found fonctionnal.
This validates the physical layer of the communication.

\subsubsection{Master-slave communication}

The complete communication is then tested.
The microcontroller is used to send command frames such as reading requests on configuration data.
The testchip is supposed to respond each time with some specific data in return.
Unfortunately, very severe issues were found with this system.
After extensive debugging and testing, it was concluded that the communication could not work as-is and needed corrections.
As will be detailed later, the issue did not originally appear in the set of simulations ran during design to validate the entire test vehicle.
It required complex debugging steps such as parasitic devices extraction from layout to only partially reproduce the issue.

% How is the test done
The test is conducted by sending a reading frame multiple times.
By design, the integrated communication system expects a \textit{clk} signal with a frequency lower than 2MHz, and an \textit{en1} signal that must be high for a single rising clock edge.
With both these criteria met, the communication system should return the data on the \textit{data} pin.

% What is expected
By design, each frame incorporates a few mechanisms to ensure the returned data is correct.
Any \textit{data} frame must start with binary code 1010 and end with binary code 01.
The \textit{en1} signal set on the external input is propagated from one read cell to the next one.
The last cell of the chain is connected to the output pin \textit{en\_out}.
Confidence in the data is increased if the \textit{enable} signal (a pulse with a width of a single clock period) is observed on \textit{en\_out},

% What are the results
After attempting multiple readings with the same board, with different boards and different clock frequencies, the results are inconsistent.
Sometimes, the communication system returns an incomplete frame like in Fig. \ref{fig:read-only-partial-frame}.
In this case, the enable signal is not propagated correctly through the chain.
The \textit{en\_out} signal (not displayed here) stays low all the time.

\begin{figure}[!h]
  \centering
  \includegraphics[width=0.8\textwidth]{src/3/figures/read_only_partial_frame.pdf}
  \caption{Read-only partially returned frame (4ms/div, 5V/div) - clock frequency 1kHz}
  \label{fig:read-only-partial-frame}
\end{figure}

In other cases, the chain returns a complete but corrupted frame like in Fig. \ref{fig:read-only-full-frame}.
The enable is correctly propagated, and is visible on \textit{en\_out}, however the data start pattern is not correct (b'1011 instead of b'1010) and some intermediate digital values are not clearly defined.

\begin{figure}[!h]
  \centering
  \includegraphics[width=0.8\textwidth]{src/3/figures/read_only_corrupted_frame.pdf}
  \caption{Read-only corrupted frame (4ms/div, 5V/div) - clock frequency 1kHz}
  \label{fig:read-only-full-frame}
\end{figure}

% Potential cause - clock freq
The clock frequency was initially suspected as a root cause for the problem.
Previous measurements were taken for a slow 1kHz clock frequency, which is far slower than the upper limit 2MHz.
Increasing the clock frequency at 10kHz, 100kHz and 1MHz did not improve the situation.

% Potential cause - delays
The impact of parasitic delays causing the communication system to malfunction was also suspected.
Multiple simulations of the entire read chain were run by placing large delays in multiple locations.
No failure could be highlighted using this method.

% Potential cause - parasitic RC
Afterwards, RC parasitic device extraction was performed from the layout.
All previous validation simulations were ran, without any success in finding the issue.
For all simulations, the read chain performed correctly.

% Potential cause - susbstrate coupling
Susbtrate couplings were also suspected.
Some reading issues could be reproduced, but this type of coupling is highly unlikely to happen for clock frequencies of 1kHz.

\subsubsection{Future work}

% Multiple bonding diagrams
A second revision of the testchip is planned.
Since the investigation on the communication system is unconclusive, it was decided to remove it until the root cause of failure could be determined.
Instead, detectors will be directly connected to pads, and most test or validation structure will be removed.
To overcome the problem of low available pin count, it is possible to put on silicon more pads than the amount of external pins.
Then, two different bonding diagrams will be made to connect each part of the available pads, to be able to use all detectors.
